\noindent\textbf{Resultado.} $y=0$ y las ramas de $xy+\ln|y|=C$ con $y\ne0$ y $xy+1\ne0$.

\noindent\textit{Desarrollo.} Escriba $M=y^2$, $N=xy+1$. No es exacta porque $M_y=2y$ y $N_x=y$. Para $y\ne0$, $(N_x-M_y)/M=-1/y$ da el factor integrante $\mu(y)=1/y$. Así $y\,dx+(x+1/y)\,dy=0$, cuya primitiva es $F=xy+\ln|y|$. La ecuación de nivel define localmente un gráfico $y(x)$ donde $F_y=x+1/y\ne0$. Su derivación da $y+(x+1/y)y^{\prime}=0$ y recupera la EDO multiplicando por $y$. En $xy+1=0$ y $y\ne0$, la ecuación original exigiría $y^2=0$, imposible; esos puntos no pertenecen a un gráfico solución clásico. La división por $y$ excluyó $y=0$, que sí satisface la EDO y debe añadirse. El campo explícito es suave cerca de cada punto de $y=0$, por lo que no se admiten empalmes arbitrarios con ramas no nulas.
